\chapter{Cadrage du Projet}
\label{chap:contexte}

\section{Introduction}
Ce chapitre présente le contexte dans lequel ce projet a été réalisé. Il introduit
\textbf{CIVCO BTP} (\textit{CivCo Building \& Engineering}) en tant qu'organisme d'accueil,
décrit la problématique métier qui a motivé le projet, et définit les objectifs ainsi que la
méthodologie adoptés tout au long du stage.

\section{Présentation de l'organisme d'accueil}

\subsection{Historique et identité}

\textbf{CivCo Building \& Engineering} est un \textbf{BET} (\textit{Bureau d'Études Techniques})
marocain, basé à \textbf{Casablanca}, spécialisé dans le génie civil, le bâtiment,
l'ingénierie des infrastructures et les solutions informatiques au service de la gestion et
du suivi des projets (Figure~\ref{fig:civco_logo}).

\begin{figure}[H]
  \centering
  
\includegraphics[width=0.35\textwidth]{civco_logo.png}
  \caption{Logo de CIVCO BTP --- CivCo Building \& Engineering}
  \label{fig:civco_logo}
\end{figure}

L'entreprise a été créée pour accompagner les maîtres d'ouvrage, les architectes et les
entreprises de construction dans la conception, l'étude et le suivi technique de leurs
projets. Au fil de son développement, elle a constitué un pôle informatique chargé de la
transformation numérique interne et du développement d'applications répondant aux besoins
des différents services.

Grâce à son expertise et au respect des normes nationales et internationales, CIVCO BTP
intervient sur des projets de tailles variées : bâtiments résidentiels, ouvrages industriels
et infrastructures publiques. Sa politique repose sur l'innovation, la qualité, la sécurité
des ouvrages et la satisfaction client, en intégrant les technologies de l'ingénierie et de
l'informatique.

\subsection{Organisation de l'entreprise}

L'organisation de CIVCO BTP repose sur une structure fonctionnelle favorisant la
collaboration entre les services techniques, administratifs et informatiques. Les principaux
services sont les suivants :

\begin{itemize}
  \item \textbf{Direction Générale} : définition de la stratégie, supervision des activités
  et développement commercial.
  \item \textbf{Service Études Techniques} : études de conception, calculs de structures
  et production des plans techniques.
  \item \textbf{Service Suivi des Travaux} : contrôle de l'exécution, respect des délais
  et suivi des chantiers.
  \item \textbf{Service Administration et Finance} : ressources humaines, comptabilité,
  achats et gestion administrative.
  \item \textbf{Service Qualité et Sécurité} : contrôle qualité, respect des normes et
  gestion de la sécurité.
  \item \textbf{Service Informatique, Réseaux et Développement} : développement
  d'applications internes et web, administration des réseaux et serveurs, gestion des
  bases de données, support utilisateurs, cybersécurité et digitalisation des processus.
\end{itemize}

La Figure~\ref{fig:civco_org} présente la répartition fonctionnelle des services au sein
de l'entreprise.

\begin{figure}[H]
  \centering
  \includegraphics[width=0.85\textwidth]{departments.png}
  \caption{Organisation fonctionnelle de CIVCO BTP}
  \label{fig:civco_org}
\end{figure}

Ce stage s'est déroulé au sein du \textbf{Service Informatique, Réseaux et Développement},
en étroite collaboration avec les équipes métier concernées par la gestion des projets.

\subsection{Domaines d'activité}

Les activités de CIVCO BTP couvrent plusieurs domaines complémentaires :

\medskip

\textbf{Génie civil et études de structures.} Études de bâtiments résidentiels et
industriels, ouvrages publics et infrastructures ; calcul en béton armé, structures
métalliques, dimensionnement des fondations et vérification de stabilité.

\medskip

\textbf{Assistance technique et conseil.} Assistance à maîtrise d'ouvrage, contrôle des
travaux, coordination des intervenants, réception des ouvrages, expertise technique,
études de faisabilité et optimisation des solutions constructives.

\medskip

\textbf{Informatique et développement.} Développement d'applications web et de logiciels
métiers, conception et administration de bases de données, maintenance applicative,
sécurisation des systèmes d'information et déploiement de solutions numériques au service
de la productivité et du suivi des projets.

\medskip

\textbf{Exemples de projets réalisés.} CIVCO BTP intervient sur des chantiers de natures
variées, illustrant la diversité de son portefeuille :

\begin{itemize}
  \item \textbf{Bâtiments résidentiels} : immeubles collectifs, villas haut standing,
  lotissements avec voirie et réseaux.
  \item \textbf{Ouvrages industriels} : halls de production, entrepôts logistiques,
  extensions d'usines existantes.
  \item \textbf{Infrastructures publiques} : voiries urbaines, ouvrages d'assainissement,
  aménagements d'espaces publics.
  \item \textbf{Rénovation et réhabilitation} : mise aux normes parasismiques, extension
  de bâtiments patrimoniaux, restructuration de locaux tertiaires.
\end{itemize}

Cette diversité impose à l'entreprise une capacité d'adaptation et un suivi rigoureux
de chaque affaire, renforçant le besoin d'un outil de gestion centralisé.

\subsection{Moyens humains et matériels}

\textbf{Moyens humains.} L'entreprise s'appuie sur une équipe pluridisciplinaire :
direction générale, ingénieurs génie civil et structures, dessinateurs-projeteurs,
conducteurs de travaux, techniciens, développeurs informatiques, administrateurs réseaux
et systèmes, personnel administratif et financier. Cette diversité de compétences permet
de répondre aux exigences techniques et organisationnelles des projets confiés à CIVCO BTP.

\medskip

\textbf{Moyens matériels.} Sur le plan technique, l'entreprise dispose d'un parc
informatique (postes de travail, serveurs, équipements réseau), de logiciels de conception
(AutoCAD, Revit, Robot Structural Analysis), d'outils de développement (Visual Studio Code,
Git, GitHub) et de systèmes de gestion de bases de données (MySQL, PostgreSQL, SQL Server).
Des véhicules de chantier, une connectivité haut débit et des plateformes collaboratives
complètent l'environnement de travail.

\section{Analyse de l'existant}

Avant de définir la solution cible, une phase d'observation et d'entretiens a été
menée auprès des différents services de CIVCO BTP. L'objectif était de comprendre
comment l'information circule aujourd'hui, quels sont les points de friction et
quels gains une application intégrée pourrait apporter.

\subsection{Processus métier actuels}

Le cycle de vie d'un projet de construction chez CIVCO BTP peut se décomposer en
plusieurs étapes successives, impliquant des acteurs distincts :

\begin{enumerate}
  \item \textbf{Prospection et relation client} : identification du besoin, échanges
  préliminaires, constitution de la fiche client (coordonnées, contacts, historique).
  \item \textbf{Étude et chiffrage} : analyse technique, estimation des coûts,
  production d'un devis détaillé avec lignes de prestations.
  \item \textbf{Contractualisation} : rédaction du contrat, éventuels avenants en cas
  de modification du périmètre ou du budget.
  \item \textbf{Planification} : découpage en phases, affectation de l'équipe,
  définition des tâches opérationnelles et des jalons.
  \item \textbf{Exécution} : suivi de chantier, mise à jour de l'avancement,
  production des bons de livraison et des documents de suivi.
  \item \textbf{Facturation et clôture} : émission des factures, suivi des paiements,
  archivage du dossier projet.
\end{enumerate}

Chaque étape génère des documents et des données qui doivent être partagés entre
le service commercial, les ingénieurs, les conducteurs de travaux, la comptabilité
et parfois le client final. Aujourd'hui, cette circulation repose en grande partie
sur des échanges manuels.

\subsection{Cartographie des outils utilisés}

Le Tableau~\ref{tab:outils_existants} synthétise les outils actuellement mobilisés
par service et leurs limites respectives.

\begin{table}[H]
\caption{Outils actuels et limites identifiées}
\label{tab:outils_existants}
\centering
\small
\begin{tabular}{|l|p{4.2cm}|p{5.8cm}|}
\hline
\textbf{Service} & \textbf{Outils} & \textbf{Limites observées} \\
\hline
Commercial & Excel, Word, courriel & Devis non centralisés, ressaisies \\
\hline
Études / Projets & AutoCAD, Revit, fichiers partagés & Pas de lien avec le suivi financier \\
\hline
Chantier & WhatsApp, tableurs & Informations éphémères, peu traçables \\
\hline
Comptabilité & Logiciel comptable, PDF & Saisie manuelle des factures \\
\hline
Direction & Rapports ponctuels & Absence de tableau de bord temps réel \\
\hline
\end{tabular}
\end{table}

Cette cartographie a confirmé l'absence d'un \textbf{référentiel unique} reliant
clients, projets, documents et équipes. Chaque service dispose de sa propre
organisation documentaire, ce qui complique la consolidation des indicateurs et
la reprise d'information lors d'un changement d'interlocuteur.

\subsection{Besoins exprimés par les services}

Les entretiens menés avec l'encadrant professionnel et les responsables de service
ont permis de formaliser les attentes prioritaires :

\begin{itemize}
  \item \textbf{Direction} : vision consolidée de l'activité (projets actifs,
  retards, charge des équipes).
  \item \textbf{Commercial} : création rapide de devis, historique client,
  conversion simplifiée en facture.
  \item \textbf{Ingénieurs / chefs de projet} : suivi des phases, des tâches et
  de l'équipe affectée à chaque chantier.
  \item \textbf{Comptabilité} : accès aux factures, suivi des paiements et état
  de facturation par projet.
  \item \textbf{Clients} : consultation de l'avancement et validation des
  documents sans dépendre des échanges par courriel.
\end{itemize}

Ces besoins ont orienté le périmètre fonctionnel de l'application développée.

\section{Contexte et problématique}

\subsection{Constats et limites des pratiques actuelles}

Au quotidien, la gestion des projets de construction chez CIVCO BTP mobilise plusieurs
services : études techniques, suivi de travaux, administration, finance et qualité.
Chaque chantier génère un volume important d'informations --- clients, phases, tâches,
devis, factures, contrats, bons de livraison --- qui doivent être partagées entre les
intervenants.

Or, une partie de ces processus repose encore sur des outils non intégrés : fichiers
Excel, documents Word ou PDF isolés, échanges par courriel ou messagerie instantanée.
Cette dispersion entraîne plusieurs difficultés :

\begin{itemize}
  \item \textbf{Fragmentation de l'information} : les données d'un même projet sont
  éparpillées entre plusieurs fichiers et supports.
  \item \textbf{Manque de visibilité} : difficulté à disposer d'une vue d'ensemble sur
  l'avancement des chantiers et des tâches en cours.
  \item \textbf{Risque d'erreurs} : ressaisies manuelles, versions multiples de documents
  et absence de source unique de vérité.
  \item \textbf{Lenteur de coordination} : dépendance aux échanges informels entre
  services pour obtenir une information à jour.
\end{itemize}

\subsection{Problématique}

Face à l'augmentation du nombre de projets et à la complexité croissante de leur suivi,
CIVCO BTP a identifié le besoin de disposer d'un \textbf{outil numérique centralisé} dédié
à la gestion et à la planification de ses projets de construction.

\medskip

La problématique de ce projet peut se formuler ainsi :

\begin{quote}
\textit{Comment concevoir et développer une application web permettant à CIVCO BTP de
centraliser, structurer et piloter l'ensemble du cycle de vie de ses projets de
construction, tout en s'adaptant aux différents profils métier de l'entreprise ?}
\end{quote}

\section{Périmètre du projet}

\subsection{Fonctionnalités incluses}

Le périmètre retenu pour ce stage couvre les modules suivants, considérés comme
prioritaires pour une première version exploitable en interne :

\begin{itemize}
  \item Authentification sécurisée et gestion des rôles (\textbf{RBAC}).
  \item Gestion des projets (phases, équipe, avancement, géolocalisation).
  \item Gestion des clients et provisionnement du portail externe.
  \item Production des documents commerciaux (devis, factures, bons de livraison).
  \item Suivi contractuel (contrats, avenants, modèles de documents).
  \item Planification et suivi des tâches opérationnelles.
  \item Tableau de bord, carte des chantiers et historique des actions.
  \item Messagerie projet entre équipe interne et client.
\end{itemize}

\subsection{Limites et hors périmètre}

Certains sujets ont été volontairement exclus ou reportés à une phase ultérieure,
afin de concentrer l'effort de développement sur le cœur métier :

\begin{itemize}
  \item \textbf{Comptabilité générale} : l'application gère les factures et
  paiements projet, mais ne remplace pas le logiciel comptable de l'entreprise.
  \item \textbf{CAO / BIM} : les outils de conception (AutoCAD, Revit) restent
  indépendants ; seules les métadonnées projet sont centralisées.
  \item \textbf{Gestion des stocks} : non couverte dans cette version.
  \item \textbf{Application mobile native} : l'interface responsive web est
  privilégiée ; une application dédiée est envisagée en perspective.
  \item \textbf{Interopérabilité ERP tiers} : pas d'intégration avec des solutions
  externes (SAP, Sage, etc.) dans le cadre de ce stage.
\end{itemize}

Cette délimitation garantit un livrable cohérent et testable dans la durée du stage,
tout en laissant une base évolutive pour les extensions futures.

\section{Objectifs du projet}

\subsection{Objectif général}

Concevoir et développer une \textbf{application web} de gestion et de planification des
projets de construction pour CIVCO BTP, capable de couvrir les principaux processus métier
liés au suivi des chantiers et à la production des documents commerciaux.

\subsection{Objectifs fonctionnels}

\begin{itemize}
  \item \textbf{Gestion des projets} : créer, consulter et suivre les chantiers par phases,
  avec affectation des équipes et suivi de l'avancement.
  \item \textbf{Gestion des clients} : centraliser les informations clients, contacts et
  historique des projets associés.
  \item \textbf{Documents commerciaux} : produire et suivre les devis, factures et bons de
  livraison (\textbf{BL}) liés aux projets.
  \item \textbf{Suivi contractuel} : gérer les contrats et les avenants associés aux
  chantiers.
  \item \textbf{Planification des tâches} : organiser le travail opérationnel par phases
  et attribuer les tâches aux collaborateurs.
  \item \textbf{Tableau de bord} : offrir une vue synthétique de l'activité (projets en
  cours, indicateurs clés, tâches prioritaires).
  \item \textbf{Contrôle d'accès} : garantir que chaque utilisateur accède uniquement aux
  fonctionnalités correspondant à son rôle (\textbf{RBAC}).
\end{itemize}

\subsection{Objectifs techniques}

\begin{itemize}
  \item Mettre en place une architecture \textit{full-stack} avec un \textit{backend} API
  \textbf{Laravel} / \textbf{PHP} et une interface \textbf{React} en \textbf{SPA}.
  \item Concevoir une base de données relationnelle (\textbf{SQL}) modélisant les entités
  métier (projets, clients, documents, utilisateurs, rôles).
  \item Exposer des services \textbf{REST} sécurisés pour la communication entre le
  \textit{frontend} et le \textit{backend}.
  \item Assurer une interface ergonomique, responsive et adaptée à un usage professionnel
  quotidien.
  \item Valider le bon fonctionnement des modules développés par des tests manuels et,
  le cas échéant, des tests automatisés.
\end{itemize}

\section{Méthodologie adoptée}

\subsection{Démarche de développement}

Le projet a été mené selon une approche \textbf{itérative et incrémentale}, adaptée au
contexte d'un stage en entreprise. Le travail s'est déroulé en plusieurs étapes :

\begin{enumerate}
  \item \textbf{Analyse des besoins} : recueil des attentes auprès de l'encadrant
  professionnel et identification des modules prioritaires.
  \item \textbf{Conception} : modélisation des entités métier, définition de l'architecture
  technique et des cas d'utilisation principaux.
  \item \textbf{Développement par modules} : implémentation progressive des fonctionnalités
  (projets, clients, documents commerciaux, tâches, tableau de bord, etc.).
  \item \textbf{Tests et validation} : vérification du comportement de chaque module et
  correction des anomalies identifiées.
  \item \textbf{Documentation} : rédaction du présent rapport et documentation technique
  du code produit.
\end{enumerate}

Chaque itération visait à livrer un incrément fonctionnel, démontrable lors des réunions
de suivi, avant de passer au module suivant.

\subsection{Gestion de projet avec Trello}

Le suivi opérationnel du stage a été organisé à l'aide de \textbf{Trello}, outil de
gestion de projet en ligne basé sur la méthode \textbf{Kanban}. Un tableau dédié au
projet regroupait les tâches sous forme de cartes réparties en colonnes selon leur
état : \textit{À faire}, \textit{En cours}, \textit{En revue} et \textit{Terminé}.

\medskip

Chaque carte correspondait à une fonctionnalité, une correction ou une tâche de
conception (par exemple : « Module devis --- API backend », « Écran tableau de bord »,
« Tests RoleAccess »). Les cartes précisaient une brève description, une priorité et,
le cas échéant, une échéance. Cette organisation a permis de :

\begin{itemize}
  \item visualiser l'avancement global du projet en un coup d'œil ;
  \item prioriser les modules à développer en accord avec l'encadrant professionnel ;
  \item conserver une trace des tâches réalisées et de celles restant à traiter ;
  \item préparer les démonstrations lors des réunions de suivi hebdomadaires.
\end{itemize}

Trello a ainsi complété Git/GitHub : le premier pour la planification et le suivi
des livrables métier, le second pour le versionnement du code source.

\subsection{Suivi académique et professionnel}

\textbf{Suivi académique (EMSI Casablanca) :}
\begin{itemize}
  \item Encadrant pédagogique : \textbf{Pr. Mohammed El Assad}
  \item Réunions régulières de suivi, validation de la démarche et préparation de la
  soutenance finale.
\end{itemize}

\textbf{Suivi professionnel (CIVCO BTP) :}
\begin{itemize}
  \item Encadrant professionnel : \textbf{M. Mohamed BENABDELLAH-CHAOUNI}
  \item Réunions de suivi pour présenter l'avancement, valider les choix fonctionnels et
  recueillir les retours métier.
  \item Alignement continu du développement avec les besoins réels de l'entreprise.
\end{itemize}

\subsection{Outils et environnement de travail}

\textbf{Outils de développement :}
\begin{itemize}
  \item \textbf{Trello} : gestion de projet et suivi des tâches (tableau Kanban).
  \item \textbf{Git / GitHub} : gestion de versions et suivi des évolutions du code.
  \item \textbf{Visual Studio Code} : environnement de développement principal.
  \item \textbf{Laravel} : framework \textit{backend} (\textbf{MVC}, \textbf{ORM} Eloquent).
  \item \textbf{React + Vite} : interface utilisateur (\textbf{SPA}).
  \item \textbf{MySQL / PostgreSQL} : base de données relationnelle.
\end{itemize}

\textbf{Environnements :}
\begin{itemize}
  \item \textbf{Développement} : poste local avec serveur Laravel et serveur de
  développement Vite (\textit{hot-reload}).
  \item \textbf{Tests} : validation manuelle des parcours utilisateurs et tests
  fonctionnels sur les modules livrés.
\end{itemize}

\section{Planification du stage}

Le stage s'est déroulé sur la période \textbf{avril -- juillet 2026}, au sein du service
informatique de CIVCO BTP. La planification générale du projet peut se résumer comme suit :

\begin{table}[H]
\centering
\caption{Phases principales du projet de stage}
\label{tab:planification}
\begin{tabularx}{\textwidth}{|l|X|}
\hline
\textbf{Phase} & \textbf{Activités principales} \\
\hline
Analyse & Étude de l'existant, recueil des besoins, définition du périmètre \\
\hline
Conception & Modélisation UML, architecture technique, maquettage des interfaces \\
\hline
Développement & Implémentation progressive des modules (\textit{backend} et \textit{frontend}) \\
\hline
Tests & Vérification fonctionnelle, corrections et stabilisation \\
\hline
Rédaction & Rédaction du rapport de fin d'études et préparation de la soutenance \\
\hline
\end{tabularx}
\end{table}

\subsection{Planning détaillé par période}

Le Tableau~\ref{tab:planning_detail} détaille la répartition des activités sur
les quatre mois du stage. Chaque période correspond à un jalon de démonstration
validé avec l'encadrant professionnel.

\begin{table}[H]
\centering
\caption{Planning détaillé du stage (avril -- juillet 2026)}
\label{tab:planning_detail}
\small
\begin{tabularx}{\textwidth}{|l|X|}
\hline
\textbf{Période} & \textbf{Livrables et activités} \\
\hline
Avril 2026 &
Analyse de l'existant, entretiens métier, rédaction du cahier des charges fonctionnel,
validation du périmètre avec l'encadrant professionnel \\
\hline
Mai 2026 &
Conception UML (cas d'utilisation, séquence, classes), modélisation de la base de données,
architecture technique three-tier, maquettage des écrans principaux \\
\hline
Juin 2026 &
Développement backend Laravel (authentification, projets, clients, devis, factures),
mise en place du frontend React (navigation, tableau de bord), premières démonstrations
des modules livrés \\
\hline
Juillet 2026 &
Portail client, tâches, configuration des rôles, tests PHPUnit et validation manuelle,
stabilisation, rédaction du rapport de fin d'études et préparation de la soutenance \\
\hline
\end{tabularx}
\end{table}

\medskip

Cette planification incrémentale a permis de livrer des modules testables au fur
et à mesure, tout en ajustant les priorités en fonction des retours terrain.
Par exemple, le module d'impression documentaire et le suivi des copies officielles
ont été intégrés en cours de développement suite à une demande exprimée par la
direction administrative.

\newpage
\section*{Conclusion}

Ce chapitre a posé le cadre du projet de fin d'études réalisé au sein de
\textbf{CIVCO BTP}. Nous avons présenté l'organisme d'accueil, son organisation
et les domaines d'activité du bureau d'études, afin de situer le projet dans son
contexte métier BTP. L'analyse de l'existant a mis en évidence la dispersion des
outils (tableurs, documents isolés, échanges informels) et la difficulté à disposer
d'une vision consolidée des chantiers, des clients et des documents commerciaux.

\medskip

À partir de cette problématique, nous avons formulé les objectifs fonctionnels et
techniques de l'application : centraliser la gestion des projets, des clients, des
devis, des factures et du suivi opérationnel, tout en garantissant un contrôle
d'accès par rôles adapté aux différents profils de l'entreprise. Le périmètre retenu
couvre les modules prioritaires identifiés avec l'encadrant professionnel, dans une
logique de livrable progressif et testable.

\medskip

La méthodologie adoptée combine une approche \textbf{agile incrémentale}, un suivi
\textbf{Kanban} via \textbf{Trello} et un versionnement \textbf{Git/GitHub}.
Le planning sur quatre mois (avril -- juillet 2026) a structuré les phases d'analyse,
de conception, de développement et de validation, avec des jalons de démonstration
réguliers. Ce cadrage constitue la base sur laquelle s'appuient les chapitres
suivants.

\medskip

Le chapitre suivant expose l'état de l'art des technologies mobilisées et justifie
les choix techniques retenus --- architecture \textit{full-stack}, \textbf{Laravel},
\textbf{React} et base de données relationnelle --- pour la réalisation de
l'application.
